% ---
% title: Two warning signals: odds
% tags: finance, odds, bayes
% ---

\begin{problem}
	\begin{context}
		Banks rarely rely on a single warning signal. A borrower may be flagged
		by a model built on account data, by a model built on market data, or by
		a loan officer. Each signal changes the bank's view of the borrower, and
		the odds form of Bayes' formula is a convenient way to combine them: each
		new signal multiplies the odds of default by its likelihood ratio.
	\end{context}
	This continues Exercise~\exref{early-warning}: 2\% of a bank's borrowers
	default within a year, and an early-warning model flags 90\% of the
	borrowers who will default and 5\% of those who will not.

	(a) Compute the probability that a flagged borrower defaults using odds:
	find the prior odds of default, the likelihood ratio of a flag, the
	posterior odds and the posterior probability.

	(b) A second model has the same rates (it flags 90\% of the defaulters and
	5\% of the non-defaulters), and its errors are independent of the first
	model's errors, given whether the borrower defaults. Both models flag a
	borrower. What is the probability that she defaults?

	(c) The first model flags a borrower, but the second does not. What is the
	probability that she defaults?
\end{problem}

\begin{slideproblem}
	2\% of borrowers default ($D$). A model flags ($A_1$) 90\% of the
	defaulters and 5\% of the non-defaulters (Exercise~\exref{early-warning}).
	(a) Using odds, find $P(D\mid A_1)$. (b) A second model with the same
	rates, whose errors are independent of the first one's given $D$ or $D^c$,
	also flags her ($A_2$). Find $P(D\mid A_1A_2)$. (c) The first model flags
	her, the second does not. Find $P(D\mid A_1A_2^c)$.
\end{slideproblem}

\begin{slides}
	\small
	\exstatement
	\pause
	\sol \textbf{Rule:} new odds $=$ old odds $\times$ likelihood ratio;
	\ odds $o$ $\leftrightarrow$ probability $o/(1+o)$.
	\pause
	\begin{itemize}
		\item Prior odds of default:
		      $\dfrac{P(D)}{P(D^c)}=\dfrac{0.02}{0.98}=\dfrac{1}{49}$
		      \quad (1 defaulter for every 49 good borrowers).
		      \pause
		\item Likelihood ratio of a flag:
		      $\dfrac{P(A_1\mid D)}{P(A_1\mid D^c)}=\dfrac{0.90}{0.05}=18$
		      \quad (a flag is 18 times as likely for a defaulter).
		      \pause
		\item (a) Posterior odds: $\dfrac{1}{49}\cdot18=\dfrac{18}{49}$,
		      \quad so
		      $P(D\mid A_1)=\dfrac{18}{18+49}=\dfrac{18}{67}\approx0.269$.
	\end{itemize}
	\pause
	Same as in Exercise~\exref{early-warning}: among the flagged, 18
	defaulters for every 49 good borrowers ($180$ vs.\ $490$).
\end{slides}

\begin{slides}
	\small
	Second model: $A_2$, same rates; given $D$ (and given $D^c$) the two flags
	are conditionally independent, so the likelihood ratios \alert{multiply}.
	\pause
	\begin{itemize}
		\item (b) Both flag: likelihood ratio $18\cdot18=324$.
		      \[ \text{odds}=\frac{1}{49}\cdot18\cdot18=\frac{324}{49},\qquad
			      P(D\mid A_1A_2)=\frac{324}{324+49}=\frac{324}{373}\approx0.869. \]
		      \pause
		\item (c) Likelihood ratio of \emph{no} flag:
		      $\dfrac{P(A_2^c\mid D)}{P(A_2^c\mid D^c)}=\dfrac{0.10}{0.95}
			      =\dfrac{2}{19}$ \quad (it lowers the odds).
		      \pause
		      \[ \text{odds}=\frac{1}{49}\cdot18\cdot\frac{2}{19}=\frac{36}{931},
			      \qquad P(D\mid A_1A_2^c)=\frac{36}{36+931}=\frac{36}{967}
			      \approx0.0372. \]
		      \pause
	\end{itemize}
	\alert{Interpretation:} two flags make default likely (87\%); one flag
	that the other model does not confirm leaves only about 4\%.
\end{slides}

\begin{solution}
	\idea Use Bayes' formula in odds form: the new odds of default are the old
	odds times the likelihood ratio of the observed signal. With two signals
	that are independent given $D$ and given $D^c$, multiply by both likelihood
	ratios. At the end, turn odds $o$ back into the probability $o/(1+o)$.

	Notation: $D$ = ``the borrower defaults'', $A_1$ = ``the first model flags
	her'', $A_2$ = ``the second model flags her''. Given: $P(D)=0.02$,
	$P(A_i\mid D)=0.90$, $P(A_i\mid D^c)=0.05$ for $i=1,2$.

	\step{The rule.} For an event $H$ and evidence $E$, Bayes' formula in odds
	form reads
	\[ \frac{P(H\mid E)}{P(H^c\mid E)}=\frac{P(H)}{P(H^c)}\cdot
		\frac{P(E\mid H)}{P(E\mid H^c)}. \]
	The first fraction on the right is the \emph{prior odds} (before the
	evidence), the second is the \emph{likelihood ratio} of the evidence: how
	many times more likely the evidence is if $H$ is true than if it is false.
	If the odds are $o$, the probability is $o/(1+o)$: odds of $o$ to $1$ mean
	$o$ ``yes'' cases for every 1 ``no'' case, out of $o+1$ cases.

	\step{(a) Prior odds.} Before any signal,
	\[ \frac{P(D)}{P(D^c)}=\frac{0.02}{0.98}=\frac{1}{49}. \]
	For every borrower who defaults, there are 49 who do not.

	\step{(a) Likelihood ratio of a flag.}
	\[ \frac{P(A_1\mid D)}{P(A_1\mid D^c)}=\frac{0.90}{0.05}=18. \]
	A flag is 18 times as likely for a borrower who will default as for one who
	will not. This is strong evidence for default.

	\step{(a) Posterior odds and probability.} Multiply:
	\[ \frac{P(D\mid A_1)}{P(D^c\mid A_1)}=\frac{1}{49}\cdot18=\frac{18}{49}. \]
	So among the flagged borrowers there are 18 defaulters for every 49 good
	borrowers, and
	\[ P(D\mid A_1)=\frac{18/49}{1+18/49}=\frac{18}{18+49}=\frac{18}{67}
		\approx0.269. \]
	This agrees with Exercise~\exref{early-warning}: of 10\,000 borrowers,
	$180$ defaulters and $490$ good borrowers are flagged, and
	$180:490=18:49$. Even though the likelihood ratio is large, the prior odds
	are so small that the posterior odds are still below 1.

	\step{(b) Two flags.} The two models' errors are independent given whether
	the borrower defaults. This means
	\[ \begin{aligned}
			P(A_1A_2\mid D)   & =P(A_1\mid D)\,P(A_2\mid D)=0.90\cdot0.90,       \\
			P(A_1A_2\mid D^c) & =P(A_1\mid D^c)\,P(A_2\mid D^c)=0.05\cdot0.05.
		\end{aligned} \]
	So the likelihood ratio of ``both flag'' is the product of the two
	likelihood ratios:
	\[ \frac{P(A_1A_2\mid D)}{P(A_1A_2\mid D^c)}=\frac{0.90}{0.05}\cdot
		\frac{0.90}{0.05}=18\cdot18=324. \]
	The posterior odds are
	\[ \frac{1}{49}\cdot324=\frac{324}{49}, \]
	and the probability is
	\[ P(D\mid A_1A_2)=\frac{324}{324+49}=\frac{324}{373}\approx0.869. \]
	Equivalently: start from the odds $18/49$ after the first flag and multiply
	by $18$ again for the second flag.

	\step{(c) A flag and no flag.} The second model does \emph{not} flag the
	borrower. This evidence has its own likelihood ratio, computed from the
	complements:
	\[ \frac{P(A_2^c\mid D)}{P(A_2^c\mid D^c)}=\frac{1-0.90}{1-0.05}
		=\frac{0.10}{0.95}=\frac{2}{19}. \]
	It is less than 1: a missing flag is more likely for a good borrower, so it
	lowers the odds of default. By the same independence (now applied to $A_1$
	and $A_2^c$, which is also allowed, because complements of conditionally
	independent events are conditionally independent too), the likelihood
	ratios multiply:
	\[ \frac{1}{49}\cdot18\cdot\frac{2}{19}=\frac{36}{931}, \]
	and
	\[ P(D\mid A_1A_2^c)=\frac{36}{36+931}=\frac{36}{967}\approx0.0372. \]
	\begin{answer}
		(a) prior odds $1/49$, likelihood ratio $18$, posterior odds $18/49$,
		$P(D\mid A_1)=18/67\approx0.269$ \quad (b) $324/373\approx0.869$ \quad
		(c) $36/967\approx0.0372$
	\end{answer}
	\begin{remark}
		The odds form makes combining signals easy: each independent signal
		multiplies the odds by its likelihood ratio, in any order. The ratios
		$18$ and $2/19$ in (c) largely cancel ($18\cdot2/19\approx1.9$), so the
		borrower ends up only a little above the prior: 3.7\% instead of 2\%.

		\textbf{Common mistake:} multiplying \emph{probabilities} instead of
		odds, or forgetting to convert the final odds back. Odds of $324/49$
		are not a probability (they are larger than 1); the probability is
		$324/373$.
	\end{remark}
\end{solution}
